%% Resumen ejecutivo (español) y Abstract (inglés). Cada uno en su página.
\clearpage
\begin{center}
  \Universidad\par\vspace{2mm}
  \Facultad\par\vspace{2mm}
  CARRERA DE \Carrera
\end{center}
\vspace{2mm}
\noindent\hspace*{\parindent}TEMA: \Tema\par
\begin{center}
  AUTOR(A): \Autor\\
  TUTOR(A): \Tutor
\end{center}
\vspace{4mm}
\begin{center}RESUMEN EJECUTIVO\end{center}

La desnutrición crónica infantil es un indicador prioritario de salud pública en Ecuador y los microdatos de la ENSANUT 2018 se utilizan sobre todo para estadísticas descriptivas. Este trabajo desarrolló y evaluó, como prototipo académico, modelos de \textit{Machine Learning} para la predicción de ese indicador, entendida como la clasificación de la condición observada en niños de 0 a 59 meses, con evaluación específica para Tungurahua. Se aplicó la metodología OSEMN a \num{18493} registros de la ENSANUT 2018 y se compararon cuatro algoritmos con partición 80/20 estratificada y selección por PR-AUC; se eligió Random Forest con umbral de decisión de 0.42. En el conjunto de prueba (n = \num{3699}) se obtuvo sensibilidad de 0.693, precisión de 0.335 y ROC-AUC de 0.677; en Tungurahua (n = 112) la sensibilidad fue de 0.686 (IC 95\,\%: 0.533–0.833). El desempeño es moderado, con sobreajuste, y no sustituye una evaluación clínica. El modelo se expuso mediante una API REST y un tablero web; 16 de 17 casos de prueba funcional se cumplieron. Se concluye que el prototipo es factible como herramienta exploratoria y que su mejora requiere validación externa.

\noindent\textbf{DESCRIPTORES:} Indicadores de salud pública; Predicción; Desnutrición crónica infantil; Aprendizaje automático; ENSANUT; OSEMN; Tungurahua

\clearpage
\begin{center}
  \Universidad\par\vspace{2mm}
  \Facultad\par\vspace{2mm}
  CARRERA DE \Carrera
\end{center}
\vspace{2mm}
\noindent\hspace*{\parindent}TOPIC:\\
\noindent\TemaIngles\par
\vspace{2mm}
\begin{center}
  AUTHOR: \Autor\\
  ADVISOR: \Tutor
\end{center}
\vspace{4mm}
\begin{center}ABSTRACT\end{center}

Chronic child malnutrition is a priority public health indicator in Ecuador, and the microdata from the 2018 ENSANUT survey are mostly used for descriptive statistics. This work developed and evaluated, as an academic prototype, machine learning models for the prediction of that indicator, understood as the classification of the condition observed in children aged 0 to 59 months, with a specific evaluation for Tungurahua. The OSEMN methodology was applied to \num{18493} ENSANUT 2018 records and four algorithms were compared using a stratified 80/20 split and PR-AUC as selection criterion; Random Forest was selected with a decision threshold of 0.42. On the test set (n = \num{3699}) the model reached a sensitivity of 0.693, precision of 0.335 and ROC-AUC of 0.677; on the Tungurahua subset (n = 112) the sensitivity was 0.686 (95\,\% CI: 0.533–0.833). Performance is moderate, with overfitting, and does not replace a clinical assessment. The model was exposed through a REST API and a web dashboard; 16 of 17 functional test cases passed. It is concluded that the prototype is feasible as an exploratory tool and that improving it requires external validation.

\noindent\textbf{KEYWORDS:} Public health indicators; Prediction; Chronic child malnutrition; Machine learning; ENSANUT; OSEMN; Tungurahua
